\chapter{Conclusion}\label{section:conclusions}

This project demonstrated that a SORT-inspired tracker can be expressed in a hardware-friendly 
form without changing the essential predict-update structure. The methodology combined background 
subtraction, bounding-box selection, constant-velocity prediction, IoU-based gating, and a compact 
Kalman-style update path. The same data stream was then used to compare software reference output 
with the custom hardware-oriented flow, giving a practical way to validate the implementation.

The results show that the tuned configuration provides a better match to the reference tracker 
than the untuned baseline, and that the final behavior depends strongly on the selected process 
and measurement noise settings. The project therefore meets its main objective of building a 
lightweight tracker that can be described cleanly in software and mapped more easily to hardware.

The main limitation is that the current design assumes a single dominant object and relies on 
foreground segmentation, so its robustness drops when the mask becomes unstable or when the object 
is briefly lost. Future work should focus on improving detector stability, adding more flexible 
association logic, and extending the architecture to handle multiple objects while retaining the 
fixed-point and control-oriented structure developed here.